% Paper 36: the leading one-loop hydrogen Lamb shift -- assembly, a
% convention erratum, and Coulomb-Sturmian Bethe logarithms.
% GeoVac Project.  Narrowed 2026-11-07 (PI-approved, group5 advisory
% pass); the cut material is preserved verbatim in
% docs/chronicles/paper36_narrowing_cut.md.

\documentclass[11pt]{article}
\usepackage{amsmath,amssymb,amsthm}
\usepackage{booktabs}
\usepackage{array}
\usepackage[margin=1in]{geometry}
\usepackage{hyperref}

\newtheorem{observation}{Observation}

\title{The Leading One-Loop Hydrogen Lamb Shift:\\
Assembly, a Convention Erratum, and Coulomb-Sturmian Bethe Logarithms}
\author{J.~Loutey (GeoVac Project)}
\date{2026}

\begin{document}
\maketitle

% ===================================================================
\begin{abstract}
% ===================================================================
We assemble the hydrogen $2S_{1/2}$--$2P_{1/2}$ Lamb shift at leading
one-loop order, $\alpha(Z\alpha)^4$, from the standard expressions
reviewed by Eides, Grotch and Shelyuto, with the electron mass in
place of the reduced mass and the Bethe logarithms tabulated by Drake
and Swainson.  The result is 1052.19~MHz, $5.65$~MHz ($0.534\%$)
below the Lundeen--Pipkin measurement, 1057.845(9)~MHz, and $5.64$~MHz
below the 2019 measurement of Bezginov \emph{et al.},
1057.8298(32)~MHz.  This is textbook arithmetic, and the residual is
not a measure of missing multi-loop physics: the one-loop
$\alpha(Z\alpha)^5$ term alone, with coefficient
$A_{50} = 4\pi(139/128 - \tfrac12\ln 2) \approx 9.29$, adds about
$+6.9$~MHz, more than the residual.  The paper records three things.
First, an erratum: an earlier assembly used the self-energy constant
$38/45$ in place of $10/9$.  The difference is exactly the Uehling
coefficient $4/15$, which had been subtracted inside the self-energy
while the Uehling shift was also added separately; restoring $10/9$
moves the Lamb shift by $+27.13$~MHz.  Second, Bethe logarithms
computed from scratch in a Coulomb-Sturmian pseudostate basis at
exponent $Z/n$: $\ln k_0(2P)$ and $\ln k_0(3D)$ reproduce the tabulated
values to $2.4\%$ and $0.07\%$ with 40 basis functions ($2P$: $1.24\%$
with 55), while the $2S$ acceleration-form value is not converged
($-0.92\%$ at 40 functions is a crossing; $+2.98\%$ at 50).  Third, a
negative result: a model two-loop self-energy sum on the Dirac
spectrum of the three-sphere grows without bound with its cutoff and
yields no finite two-loop coefficient without renormalization
counterterms, which it does not contain.  Evaluating hydrogen Bethe
logarithms through the dynamical symmetry of the Coulomb problem is
prior art (Lieber, 1968; Huff, 1969).
\end{abstract}

% ===================================================================
\section{Introduction}
\label{sec:intro}
% ===================================================================

The $2S_{1/2}$--$2P_{1/2}$ splitting of hydrogen, first measured by
Lamb and Retherford in 1947 and explained at leading order by Bethe's
non-relativistic calculation the same year~\cite{bethe1947}, is the
canonical test of bound-state QED.  The two direct measurements used
here are Lundeen and Pipkin's, $1057.845(9)$~MHz~\cite{lundeen_pipkin1981},
and the frequency-offset separated-oscillatory-field measurement of
Bezginov \emph{et al.}, $1057.8298(32)$~MHz, obtained from a measured
$2S_{1/2}(F{=}0)$--$2P_{1/2}(F{=}1)$ interval of $909.8717(32)$~MHz
after removing the hyperfine contribution~\cite{bezginov2019}.  Both are
intervals between hyperfine centroids.  The modern theory of the
interval, order by order, is reviewed by Eides, Grotch and
Shelyuto~\cite{eides2001}; current work on the two-loop self-energy at
low $Z$ is that of Yerokhin, Harman and
Keitel~\cite{yerokhin2024_two_loop_se}.

\paragraph*{Scope of the claim.}
The $1052.19$~MHz headline is the leading one-loop assembly with
tabulated Bethe logarithms.  It is not a GeoVac computation and does
not compete with precision theory.  It omits the one-loop
$\alpha(Z\alpha)^5$ term, the reduced-mass factors $(m_r/m)^3$ of the
leading terms, recoil, the finite nuclear size, and all two-loop terms;
the residual is therefore not attributed (Sec.~\ref{sec:residual}).
Nothing in the number uses the three-sphere: no Fock projection, no
Dirac spectrum on $S^3$, no GeoVac label set.  What the paper adds is
(i)~the erratum of Sec.~\ref{sec:eides}, (ii)~the Bethe logarithms of
Sec.~\ref{sec:bethe}, computed from scratch and checked against the
Drake--Swainson table~\cite{drake_swainson}, and (iii)~the two-loop
negative of Sec.~\ref{sec:twoloop}.

\paragraph*{Prior art.}
The use of the hydrogen atom's dynamical symmetry to evaluate Bethe
logarithms is due to Lieber, who used the O(4) symmetry together with
Green's-function methods to obtain a rapidly convergent series and
values for $1S$, $2S$ and $2P$~\cite{lieber1968}, and to Huff, who
recast the calculation in the algebra of SO(4,2)~\cite{huff1969}.  The
Coulomb-Sturmian pseudostate sums of Sec.~\ref{sec:bethe} are a finite
basis version of the same spectral sums; we make no novelty claim for
the method.

\paragraph*{Backing.}
\sloppy
The assembly and the erratum are recomputed in exact and closed-form
arithmetic by the test file
\texttt{tests/test\_paper36\_lamb\_chain.py}, which also recomputes
(slow-marked) the Bethe logarithms from the Sturmian code in
\texttt{tests/paper36\_lamb\_support/}.  The $\alpha(Z\alpha)^5$
estimate is checked by \texttt{tests/test\_paper36\_a50\_term.py}.  The
two-loop model is \texttt{geovac/two\_loop\_self\_energy.py}.\par
\fussy

% ===================================================================
\section{The one-loop expressions}
\label{sec:oneloop}
% ===================================================================

At order $\alpha(Z\alpha)^4$ and with $m_r = m$, the self-energy and
Uehling vacuum-polarization shifts of the $n=2$ levels are, in Hartree
($\alpha(Z\alpha)^4 m c^2 = \alpha^3 Z^4\,\mathrm{Ha}$),
\begin{align}
\label{eq:eides_se}
\Delta E_\text{SE}(nS_{1/2}) &= \frac{\alpha^3 Z^4}{\pi n^3}
\Bigl[\tfrac{4}{3} \ln\frac{1}{(Z\alpha)^2}
- \tfrac{4}{3} \ln k_0(n, 0) + \tfrac{10}{9}\Bigr], \\
\label{eq:eides_se_2p}
\Delta E_\text{SE}(2P_{1/2}) &= \frac{\alpha^3 Z^4}{8\pi}
\Bigl[- \tfrac{4}{3} \ln k_0(2, 1) - \tfrac{1}{6}\Bigr], \\
\label{eq:uehling_vp}
\Delta E_\text{VP}(nS) &= -\frac{4}{15}\,\frac{\alpha^3 Z^4}{\pi n^3},
\end{align}
with $\Delta E_\text{VP}(2P) = 0$ at this order.  These are the entries
of Ref.~\cite{eides2001}, Sec.~VIII.A, Eq.~(40) and Table~II, rewritten
with the common factor $\alpha(Z\alpha)^4 m/(\pi n^3)$: the constant
$10/9 = 11/18 + 1/2$ is the Dirac form-factor term ($4 \times 11/72$)
plus the Pauli form-factor term ($4 \times 1/8$); the $2P_{1/2}$
constant $-1/6$ is the Pauli term $4(j(j+1) - l(l+1) - 3/4)/
(8l(l+1)(2l+1))$ at $l = 1$, $j = 1/2$; and $-4/15 = 4 \times (-1/15)$
is the Uehling term.  The Lamb shift is
$L = [\Delta E_\text{SE} + \Delta E_\text{VP}](2S) -
\Delta E_\text{SE}(2P_{1/2})$.

% ===================================================================
\section{The erratum: a double-counted Uehling term}
\label{sec:eides}
% ===================================================================

An earlier implementation of Eq.~\eqref{eq:eides_se} used the constant
$38/45$ in place of $10/9$, and also added the Uehling shift
\eqref{eq:uehling_vp}.  The difference of the two constants is
\begin{equation}
\label{eq:ls1_constant_difference}
\frac{10}{9} - \frac{38}{45} \;=\; \frac{50 - 38}{45} \;=\; \frac{4}{15},
\end{equation}
exactly the Uehling coefficient.  The earlier constant had the Uehling
term already subtracted inside the self-energy bracket, so with the
separate Uehling shift the vacuum polarization was counted twice.
Restoring $10/9$ moves the Lamb shift by
\begin{equation}
\label{eq:shift_value}
\Delta L = +\frac{4}{15} \cdot \frac{\alpha^3 Z^4}{\pi n^3}\,\mathrm{Ha}
\;=\; +27.13~\text{MHz at } n = 2,\ Z = 1.
\end{equation}

\begin{observation}[Uehling double counting]
\label{obs:uehling_double}
With the self-energy constant $38/45$ and a separate Uehling term the
leading one-loop Lamb shift is $1025.06$~MHz; with the canonical
$10/9$ it is $1052.19$~MHz.  The shift is exactly
$(4/15)\,\alpha^3 Z^4/(\pi n^3)$~Ha $= +27.13$~MHz.
\end{observation}

% ===================================================================
\section{Result}
\label{sec:result}
% ===================================================================

With the Drake--Swainson values $\ln k_0(2S) = 2.8117698931$ and
$\ln k_0(2P) = -0.0300167089$~\cite{drake_swainson}, $\alpha^{-1} =
137.035999084$ and $1~\mathrm{Ha} = 6\,579\,683\,920.502$~MHz:

\begin{center}
\begin{tabular}{l r r r}
\toprule
Quantity & earlier (MHz) & corrected (MHz) & Shift (MHz) \\
\midrule
SE $2S_{1/2}$ constant & $+38/45$ & $+10/9$ & $+4/15$ \\
SE $2S_{1/2}$ total & $+1039.31$ & $+1066.44$ & $+27.13$ \\
VP $2S_{1/2}$ (Uehling) & $-27.13$ & $-27.13$ & $0$ \\
Total $2S_{1/2}$ & $+1012.18$ & $+1039.31$ & $+27.13$ \\
SE $2P_{1/2}$ & $-12.88$ & $-12.88$ & $0$ \\
VP $2P_{1/2}$ & $0$ & $0$ & $0$ \\
\midrule
\textbf{Lamb shift} & \textbf{$+1025.06$} & \textbf{$+1052.19$} & \textbf{$+27.13$} \\
\midrule
Lundeen--Pipkin~\cite{lundeen_pipkin1981} & \multicolumn{2}{c}{$+1057.845(9)$} & \\
Residual & $-32.78$ ($-3.10\%$) & $-5.65$ ($-0.534\%$) & \\
\bottomrule
\end{tabular}
\end{center}

\noindent Against Bezginov \emph{et al.}~\cite{bezginov2019}
($1057.8298(32)$~MHz) the corrected residual is $-5.64$~MHz.

\subsection{The residual}
\label{sec:residual}

The residual is not attributed.  The assembly is the leading
$\alpha(Z\alpha)^4$ order only, at $m_r = m$.  The one-loop
$\alpha(Z\alpha)^5$ term,
\begin{equation}
\label{eq:a50}
\Delta E(nS) = \frac{\alpha(Z\alpha)^5}{\pi n^3}\, m c^2\, A_{50},
\qquad A_{50} = 4\pi\Bigl(\frac{139}{128} - \frac12\ln 2\Bigr) \approx 9.29,
\end{equation}
contributes about $+6.9$~MHz to the $2S$ level and by itself exceeds
the residual; Eides \emph{et al.}\ give $+7.13$~MHz for the complete
one-loop $\alpha(Z\alpha)^5$ correction to $2S$, including the
polarization insertions~\cite{eides2001} (Table~III).  Recoil, the
finite nuclear size, and two-loop terms enter besides, with either
sign; the two-loop $\alpha^2(Z\alpha)^4$ entries of the same review sum to
$+0.127$~MHz on $2S$ (Table~II).  An order-by-order budget is owed
before any attribution.  An earlier decomposition of the residual into
multi-loop and non-loop pieces is withdrawn.

% ===================================================================
\section{Bethe logarithms in a Coulomb-Sturmian basis}
\label{sec:bethe}
% ===================================================================

\subsection{Definition and normalization}

For a hydrogenic state $|nlm\rangle$ (atomic units),
\begin{equation}
\label{eq:bethe_log_def}
\ln k_0(n, l) = \frac{\sum_{k} |\langle nlm | \vec{p} | k \rangle|^2
(E_k - E_{n}) \ln \bigl(|E_k - E_{n}| / \mathrm{Ry}\bigr)}
{2 Z^4 / n^3},
\end{equation}
where the sum runs over all intermediate states, bound and continuum,
and their magnetic substates.  The denominator is the same for every
$l$.  For $s$-states it equals the sum rule
$I_v(nS) = \sum_k |\langle nS | \vec{p} | k \rangle|^2 (E_k - E_n) =
\tfrac12\langle \nabla^2 V \rangle = 2Z^4/n^3$; for $l > 0$ the
corresponding sum vanishes, because $\psi_{nl}(0) = 0$, while the
numerator of Eq.~\eqref{eq:bethe_log_def} stays finite.  This uniform
normalization is the one in which Eq.~\eqref{eq:eides_se_2p} and the
Drake--Swainson table are written.

In the code the numerator is summed over the initial magnetic
substates as well, which multiplies it by $2l + 1$; the code therefore
divides by $D(n,l) = 2(2l+1)Z^4/n^3$.  The factor $2l + 1$ is this
bookkeeping and nothing more: without it the computed $2P$ logarithm
would be three times the tabulated one.

\subsection{Basis}

The radial basis is
$\chi_{k,l}(r) = r^{l+1} e^{-\lambda r} L_k^{(2l+1)}(2 \lambda r)$,
$k = 0, \ldots, N - 1$, at $\lambda = Z/n$ for target level $n$, which
contains the exact $n$-th hydrogenic radial function.  All overlap,
kinetic, potential, dipole and $1/r^2$ matrix elements are exact
polynomial-times-exponential integrals evaluated at 50 digits; the
Hamiltonian is diagonalized in the basis after canonical
orthogonalization, and the pseudostates of the $l \pm 1$ channels
discretize the intermediate spectrum.

\subsection{\texorpdfstring{Velocity form for $l > 0$}{Velocity form for l > 0}}

With $|\langle nl|\vec p|k\rangle|^2 = (E_k - E_n)^2\, l_>\,
|\langle nl|r|k\rangle|^2$ (summed over substates) the numerator of
Eq.~\eqref{eq:bethe_log_def} converges for $l > 0$ and is evaluated
directly.  Against the Drake--Swainson values~\cite{drake_swainson}:
\begin{itemize}
\setlength\itemsep{0pt}
\item $\ln k_0(2P) = -0.03074$ at $N = 40$, $+2.40\%$ from
$-0.0300167$; $-0.03039$ at $N = 55$, $+1.24\%$.
\item $\ln k_0(3D) = -0.005236$ at $N = 40$, $+0.07\%$ from
$-0.0052321$.
\end{itemize}

\subsection{\texorpdfstring{Acceleration form for $s$-states}{Acceleration form for s-states}}

For $s$-states the velocity form converges slowly.  The acceleration
form replaces $\vec p$ by $\nabla V = Z\hat r/r^2$,
\begin{equation}
\ln k_0 = \frac{\sum_k |\langle nl|\nabla V|k\rangle|^2
\ln\bigl(|E_k - E_n|/\mathrm{Ry}\bigr)/(E_k - E_n)}
{\sum_k |\langle nl|\nabla V|k\rangle|^2/(E_k - E_n)}.
\end{equation}
In this basis it does not converge:
\begin{itemize}
\setlength\itemsep{0pt}
\item $\ln k_0(2S) = 2.7858$ at $N = 40$ ($-0.92\%$ from $2.8117699$)
and $2.8955$ at $N = 50$ ($+2.98\%$); the $N = 40$ agreement is a
crossing.
\item $\ln k_0(1S) = 3.0019$ at $N = 20$ ($+0.60\%$ from $2.9841286$)
and $3.268$ at $N = 40$ ($+9.5\%$); again a crossing.
\end{itemize}
A converged from-scratch $2S$ logarithm is therefore not available
here, and the assembly of Sec.~\ref{sec:result} uses the tabulated
value.

% ===================================================================
\section{A two-loop negative}
\label{sec:twoloop}
% ===================================================================

A model of the two-loop self-energy was built on the Dirac spectrum
of the unit three-sphere~\cite{camporesi_higuchi}, $|\lambda_n| = n +
3/2$ with degeneracy $g_n = 2(n+1)(n+2)$, and transverse photon modes
of eigenvalue $\mu_q = q(q+2)$ and degeneracy $d_q = q(q+2)$.  For the
rainbow and crossed topologies it sums
\begin{equation}
\Sigma = \sum_{n_1,n_2,n_3,\,q_a,q_b}
\frac{\bigl(\prod_{i=1}^{4} W_i\bigr)\, g_{n_1} g_{n_2} g_{n_3}\,
d_{q_a} d_{q_b}}
{\lambda_{n_1}^2 \lambda_{n_2}^2 \lambda_{n_3}^2\, \mu_{q_a} \mu_{q_b}}
\end{equation}
over electron levels $0, \ldots, n_\text{max}$ and photon levels
allowed at each vertex, where $W_i \in \{0, 2\}$ is an SO(4)
channel count with the parity rule $n_a + n_b + q$ odd, and the
external level is $n_\text{ext} = 1$.  No bound-state wavefunction
enters; the external level stands in for hydrogen $2S$.  The
dimensionless bracket is $C = \Sigma_\text{total}/(4\pi)^2$.

\begin{table}[ht]
\caption{The model two-loop sum versus cutoff, with the channel count
as corrected on 2026-11-07 (an earlier version of the count paired the
spinor content of opposite signs of the Dirac operator, and every
entry of the table printed before that date was computed with it).}
\label{tab:ls8a_native}
\centering
\begin{tabular}{c r r r r}
\toprule
$n_\text{max}$ & rainbow & crossed & total & $C$ \\
\midrule
2 & 587.4 & 467.0 & 1054.4 & 6.68 \\
3 & 1680.9 & 1313.8 & 2994.8 & 18.96 \\
4 & 3653.0 & 2786.3 & 6439.3 & 40.78 \\
5 & 6761.9 & 5140.6 & 11902.5 & 75.37 \\
6 & 11265.7 & 8507.5 & 19773.2 & 125.22 \\
\bottomrule
\end{tabular}
\end{table}

The sum grows without bound: between $n_\text{max} = 5$ and $6$ the
local log-log slope of the total is $2.8$ (it was $3.3$ under the
earlier count, whose totals were $44.8$, $266.3$, $761.3$, $1657.8$ and
$3052.8$).  A finite two-loop coefficient requires the
wave-function and mass renormalization counterterms, $Z_2$ and $\delta
m$, which the model does not contain.  The negative is about this
model: it does not show that no spectral construction on $S^3$ can
produce the two-loop coefficient, and it makes no comparison with the
literature value, since the model has no finite value to compare.

% ===================================================================
\section{Conclusion}
\label{sec:conclusion}
% ===================================================================

The leading one-loop assembly of the hydrogen Lamb shift with tabulated
Bethe logarithms gives $1052.19$~MHz, $0.534\%$ below the
Lundeen--Pipkin measurement.  The number is standard one-loop QED with
literature inputs, and the residual is not attributed: the one-loop
$\alpha(Z\alpha)^5$ term alone is about $+6.9$~MHz.  An earlier
version of the assembly double-counted the Uehling term, and the
correction is exactly $(4/15)\,\alpha^3 Z^4/(\pi n^3)$~Ha
$= +27.13$~MHz.  Bethe logarithms computed from scratch in a
Coulomb-Sturmian pseudostate basis reproduce the tabulated $2P$ and
$3D$ values to $2.4\%$ and $0.07\%$ at 40 functions, while the $2S$
acceleration-form value has not converged.  A model two-loop
self-energy sum on the three-sphere Dirac spectrum diverges with its
cutoff and gives no finite coefficient without renormalization
counterterms.  The muonic-hydrogen transfer that an earlier version of
this paper reported is in Paper~34.

% ===================================================================
\begin{thebibliography}{99}
% ===================================================================

\bibitem{bethe1947}
H.~A.~Bethe, ``The Electromagnetic Shift of Energy Levels,''
\textit{Phys.\ Rev.}\ \textbf{72}, 339 (1947).

\bibitem{lundeen_pipkin1981}
S.~R.~Lundeen and F.~M.~Pipkin, ``Measurement of the Lamb Shift in
Hydrogen, $n = 2$,'' \textit{Phys.\ Rev.\ Lett.}\ \textbf{46}, 232--235
(1981).

\bibitem{bezginov2019}
N.~Bezginov \textit{et al.}, ``A measurement of the atomic hydrogen
Lamb shift and the proton charge radius,'' \textit{Science}\
\textbf{365}, 1007--1012 (2019), doi:10.1126/science.aau7807.

\bibitem{eides2001}
M.~I.~Eides, H.~Grotch, and V.~A.~Shelyuto, ``Theory of light
hydrogenlike atoms,'' \textit{Phys.\ Rep.}\ \textbf{342}, 63--261
(2001); arXiv:hep-ph/0002158 (section, equation and table numbers as
in the arXiv version).

\bibitem{drake_swainson}
G.~W.~F.~Drake and R.~A.~Swainson, ``Bethe logarithms for hydrogen
up to $n=20$, and approximations for two-electron atoms,''
\textit{Phys.\ Rev.\ A}\ \textbf{41}, 1243--1246 (1990).

\bibitem{lieber1968}
M.~Lieber, ``O(4) Symmetry of the Hydrogen Atom and the Lamb Shift,''
\textit{Phys.\ Rev.}\ \textbf{174}, 2037--2054 (1968).

\bibitem{huff1969}
R.~W.~Huff, ``Simplified Calculation of Lamb Shift Using Algebraic
Techniques,'' \textit{Phys.\ Rev.}\ \textbf{186}, 1367--1379 (1969).

\bibitem{yerokhin2024_two_loop_se}
V.~A.~Yerokhin, Z.~Harman, and C.~H.~Keitel,
``Two-loop electron self-energy for low nuclear charges,''
\textit{Phys.\ Rev.\ Lett.}\ \textbf{133}, 251803 (2024); arXiv:2411.12459.

\bibitem{camporesi_higuchi}
R.~Camporesi and A.~Higuchi, ``On the eigenfunctions of the Dirac
operator on spheres and real hyperbolic spaces,''
\textit{J.\ Geom.\ Phys.}\ \textbf{20}, 1 (1996).

\end{thebibliography}

\end{document}
